\begin{helper}
Fix a terminal coordinate set \(U\), a blue support \(B\), a red label
\(J\), a finite branch-label type \(\mathcal T\), and a finite transcript
family \(\mathcal P\).  Let
\[
\mathsf{Realized}(\beta,\tau)
\]
mean that \((\beta,\tau)\) is a nonempty fixed \((B,J)\) red cell, where
\(\tau=(P_\tau,A_\tau,Z_\tau)\).  Assume \(B\) is a blue \(u_B\)-subset of
\(U\).  For every realized cell, assume the decoded red support
\(R_{\beta,J}\) is a red \(u_R\)-subset of \(U\), and that
\[
P_\tau,A_\tau\subseteq U,\qquad P_\tau\cap A_\tau=\emptyset,
\]
\[
B\cup R_{\beta,J}\subseteq P_\tau,\qquad Z_\tau\subseteq A_\tau,\qquad
\max(0,a-u_R-u_B)\le |Z_\tau|.
\]

Assume that the residual product tree for these fixed supports has a finite
leaf set \(\Lambda\), a leaf mass \(\nu:\Lambda\to\mathbb R\), and a leaf
assignment
\[
\iota(\beta,\tau)\in\Lambda
\]
for every formal branch/transcript pair.  The leaf masses are nonnegative
and normalized:
\[
0\le \nu(\lambda)\quad(\lambda\in\Lambda),\qquad
\sum_{\lambda\in\Lambda}\nu(\lambda)\le1.
\]
Assume also that the leaf assignment is injective on realized cells: if
\(\tau,\tau'\in\mathcal P\), both cells are realized, and
\[
\iota(\beta,\tau)=\iota(\beta',\tau'),
\]
then \(\beta=\beta'\) and \(\tau=\tau'\).  Then there are cell-indexed
residual masses \(r_{\beta,\tau}\) such that
\[
\sum_{\beta\in\mathcal T}
  \sum_{\tau\in\mathcal P:\mathsf{Realized}(\beta,\tau)}
    r_{\beta,\tau}\le1,
\]
\[
0\le r_{\beta,\tau}
\]
on realized cells, and
\[
r_{\beta,\tau}=\nu(\iota(\beta,\tau))
\]
on realized cells.
\end{helper}
\begin{proof}
Define
\[
r_{\beta,\tau}=\nu(\iota(\beta,\tau))
\]
for all formal pairs \((\beta,\tau)\).  On a realized pair this immediately
gives the displayed identity, and nonnegativity follows from the
nonnegativity of the corresponding residual leaf mass.

It remains to prove the total bound.  Let
\[
\mathcal C=\{(\beta,\tau):\tau\in\mathcal P,\,
  \mathsf{Realized}(\beta,\tau)\}
\]
be the finite set of realized cells.  The injectivity hypothesis says exactly
that the map \((\beta,\tau)\mapsto\iota(\beta,\tau)\) is injective on
\(\mathcal C\).  Therefore the sum of the cell residual masses is the sum of
the same leaf masses over the image \(\iota(\mathcal C)\), with no leaf
counted twice:
\[
\sum_{(\beta,\tau)\in\mathcal C} r_{\beta,\tau}
=
\sum_{\lambda\in\iota(\mathcal C)} \nu(\lambda).
\]
The image \(\iota(\mathcal C)\) is a subset of the finite residual leaf set
\(\Lambda\).  Since every \(\nu(\lambda)\) is nonnegative, enlarging the sum
from this image to all leaves can only increase it:
\[
\sum_{\lambda\in\iota(\mathcal C)} \nu(\lambda)
\le
\sum_{\lambda\in\Lambda}\nu(\lambda).
\]
The residual product-tree normalization gives the final upper bound
\[
\sum_{\lambda\in\Lambda}\nu(\lambda)\le1.
\]
Combining the three displayed relations proves the claimed subprobability
bound for the cell-indexed residual masses.
\end{proof}
